\chapter{The APL Source Semantics: Iota, Omega, and Reduction}
\label{ch:apl-sem}

\section{Why APL Is the User Surface}
\label{sec:apl-why}

APL is the language the numerical user actually writes. Its array-first,
symbol-dense notation is compact but its semantics are precise and
well-suited to formal treatment. MATHLIB5 accepts a deliberately small
fragment: rank-1 (vector) and rank-2 (matrix) arrays, plus the primitives
$\iota$ (iota), $\omega$ (omega), reduction ($/$), and outer product
($\circ.\times$). This chapter fixes their semantics before the IR
translation of \cref{ch:ir}.

\section{Arrays and Indexing}
\label{sec:apl-array}

An APL vector $v$ of length $n$ is indexed $v[1], v[2], \dots, v[n]$
(1-based). A matrix $M$ of shape $m\times n$ has entries $M[i;j]$ with
$1\le i\le m$, $1\le j\le n$. The translation to C$--$ (\cref{ch:cmm})
shifts to 0-based indexing; the shift is itself verified by the Layer~2
$\to$ IR contract so that no off-by-one can survive (\cref{ch:liquid}).

\section{Iota}
\label{sec:apl-iota}

The monadic iota $\iota n$ yields the vector
\[
\iota n \;=\; (1, 2, \dots, n).
\]
The kernel computes it as the integer range; in Liquid Haskell the
postcondition is
\[
\forall k\in\{1,\dots,n\}.\;\; (\iota n)_k = k,
\]
which the refinement solver discharges by induction.

\section{Omega}
\label{sec:apl-omega}

$\omega$ is the right argument of a defined function; as a standalone
glyph in our fragment it denotes the input vector to the summation. Formally,
a summation expression $\Sigma\,\omega$ has denotation
\[
\denote{\Sigma\,\omega}(v) \;=\; \sum_{k=1}^{|v|} v_k.
\]
The closed form established in \cref{ch:lean} is the identity
\[
\sum_{k=1}^{n} v_k \;=\; \frac{n}{2}\,(v_1 + v_n)
\quad\text{when } v_k \text{ is arithmetic}.
\]

\section{Reduction}
\label{sec:apl-red}

The reduction $f/$ over a vector applies the dyadic function $f$
between consecutive elements:
\[
+/\, v \;=\; v_1 + (v_2 + (\cdots + v_n)\cdots).
\]
Because $+$ is associative, the grouping is irrelevant to the value, which
lets the C$--$ code choose an efficient left-to-right fold without
changing the certified result. The Lean proof of associativity of integer
addition underpins this freedom.

\section{Outer Product}
\label{sec:apl-outer}

The outer product $f\circ.\times$ of two vectors $a$ (length $m$) and
$b$ (length $n$) produces the $m\times n$ matrix
\[
(a\circ.\times\, b)_{ij} = a_i\; f\; b_j.
\]
For $f = \times$, this is the standard outer product. The kernel's
$\Sigma$ computation uses it to build the weight matrix before
contracting rows. The IR representation (\cref{ch:ir}) stores it as a
two-dimensional array with explicit shape metadata so the later layers
never lose the $m\times n$ structure.

\section{Worked Trace}
\label{sec:apl-trace}

Consider the input vector $v = (1,2,3,4)$. The fragment
\[
\Sigma\,\omega \quad\text{on}\quad v
\]
reduces as follows in the kernel:
\begin{align*}
+/\, v &= 1+2+3+4 = 10,\\
\text{closed form} &: \frac{4}{2}(1+4) = 2\cdot 5 = 10.
\end{align*}
The two computations agree, and the Lean proof certifies that the closed
form equals the reduction for every arithmetic $v$. The artifact
propagates to the IR as the tuple $(v,\, 10,\,\text{proof-ref})$.

\section{Translation Contract to the IR}
\label{sec:apl-contract}

Layer~1 $\to$ Layer~2 is summarized by the contract
\[
\forall e\in\text{APL}_{\text{frag}}.\;\;
\denote{\text{IR}(e)} = \denote{e},
\]
where $\text{IR}(e)$ is the decorated intermediate representation. The
decoration carries the shape, the element type (integer vs.\ float), and a
reference to the closed-form certificate when one exists. This contract is
the first boundary of the pipeline and the first receipt emitted
(\cref{ch:worm}).
